\documentclass[11pt,a4paper]{article}
\input{../template/utm-lab}

\newcommand{\labnumber}{13}
\newcommand{\labtitle}{Semaphores in Your Own Kernel}
\newcommand{\studentname}{Daniel Nicorici}
\newcommand{\studentgroup}{FAF-243}

\begin{document}

\utmtitlepage
\tableofcontents
\newpage

\plainsection{Purpose of the work}

In Labs 11 and 12 semaphores were used from the C library. The purpose of this lab is to build
them inside the xv6 kernel, as three new system calls, \code{sem\_init}, \code{sem\_down} and
\code{sem\_up} (numbers 24, 25 and 26), and to see that a semaphore is only a number, a
spinlock that protects it, and the kernel's sleep and wakeup.

The work was done in my \code{\textasciitilde/xv6-riscv} tree (\code{mit-pdos/xv6-riscv}, commit
\code{06aad25}), which already contains \code{trace} from Lab 10 as call 23, cross-compiled with
\code{riscv64-linux-gnu-gcc} 16.2.1 and run on QEMU 11.1.1 with 3 CPUs.

\section{Part 1. The numbers and the user side}

The same recipe as \code{trace} in Lab 10, three times. Semaphores are numbered 0 to 7 and are
shared by all processes.

\begin{lstlisting}[caption={The numbers, the declarations, the stubs and the UPROGS lines}]
$ grep -n "sem_" kernel/syscall.h user/user.h user/usys.pl
kernel/syscall.h:25:#define SYS_sem_init 24
kernel/syscall.h:26:#define SYS_sem_down 25
kernel/syscall.h:27:#define SYS_sem_up   26
user/user.h:29:int sem_init(int, int);
user/user.h:30:int sem_down(int);
user/user.h:31:int sem_up(int);
user/usys.pl:47:entry("sem_init");
user/usys.pl:48:entry("sem_down");
user/usys.pl:49:entry("sem_up");
$ grep -n "_semtest\|_mix" Makefile
151:	$U/_semtest\
152:	$U/_mix\
\end{lstlisting}

\code{semtest} has two tests. In the first, semaphore 0 starts at 0, so the child must wait for
the parent. In the second, ping and pong take turns with two semaphores that pass the turn to
each other.

\begin{lstlisting}[style=ccode,caption={user/semtest.c}]
#include "kernel/types.h"
#include "user/user.h"

int
main(void)
{
  sem_init(0, 0);
  if (fork() == 0) {
    sem_down(0);
    printf("child: I ran after the parent\n");
    exit(0);
  }
  pause(10);
  printf("parent: done, I let the child go\n");
  sem_up(0);
  wait(0);

  sem_init(1, 1);
  sem_init(2, 0);
  if (fork() == 0) {
    for (int i = 1; i <= 3; i++) {
      sem_down(2);
      printf("  pong %d\n", i);
      sem_up(1);
    }
    exit(0);
  }
  for (int i = 1; i <= 3; i++) {
    sem_down(1);
    printf("ping %d\n", i);
    sem_up(2);
  }
  wait(0);
  printf("semtest: OK\n");
  exit(0);
}
\end{lstlisting}

\begin{lstlisting}[style=ccode,caption={user/mix.c}]
#include "kernel/types.h"
#include "user/user.h"

void
lines(char *s, int uselock)
{
  for (int k = 0; k < 5; k++) {
    if (uselock)
      sem_down(3);
    for (char *p = s; *p; p++)
      write(1, p, 1);
    if (uselock)
      sem_up(3);
  }
}

int
main(int argc, char *argv[])
{
  int uselock = (argc > 1);
  sem_init(3, 1);
  if (fork() == 0) {
    lines("bbbbbbbbbbbbbbbbbbbbbbbbbbbbbb\n", uselock);
    exit(0);
  }
  lines("AAAAAAAAAAAAAAAAAAAAAAAAAAAAAA\n", uselock);
  wait(0);
  exit(0);
}
\end{lstlisting}

\begin{lstlisting}[caption={Check 1: the calls reach a kernel that does not know them yet}]
$ semtest
3 semtest: unknown sys call 24
4 semtest: unknown sys call 25
child: I ran after the parent
parent: done, I let the child go
3 semtest: unknown sys call 26
3 semtest: unknown sys call 24
3 semtest: unknown sys call 24
3 semtest: unknown sys call 25
[... ping and pong mixed with the error lines ...]
semtest: OK
\end{lstlisting}

The line of the child comes before the line of the parent: without semaphores the child ran
first. That is what Part 2 fixes.

\section{Part 2. The kernel side}

\code{sem\_down} is the heart. While the value is 0, the process registers the address of the
semaphore as the thing it waits for, lets go of the spinlock, and sleeps. \code{sem\_up} adds 1
and wakes up everyone who sleeps on that address. This kernel sleeps in two steps,
\code{sleep\_prepare(chan)} and then \code{sleep()}.

\begin{lstlisting}[style=ccode,caption={kernel/sem.c}]
#include "types.h"
#include "riscv.h"
#include "defs.h"
#include "param.h"
#include "spinlock.h"
#include "proc.h"

#define NSEM 8

struct sem {
  struct spinlock lock;
  int value;
};

static struct sem sems[NSEM];

uint64
sys_sem_init(void)
{
  int id, value;
  argint(0, &id);
  argint(1, &value);
  if (id < 0 || id >= NSEM || value < 0)
    return -1;
  acquire(&sems[id].lock);
  sems[id].value = value;
  release(&sems[id].lock);
  return 0;
}

uint64
sys_sem_down(void)
{
  int id;
  argint(0, &id);
  if (id < 0 || id >= NSEM)
    return -1;
  struct sem *s = &sems[id];

  acquire(&s->lock);
  while (s->value == 0) {
    sleep_prepare(s);
    release(&s->lock);
    sleep();
    acquire(&s->lock);
    if (killed(myproc())) {
      release(&s->lock);
      return -1;
    }
  }
  s->value--;
  release(&s->lock);
  return 0;
}

uint64
sys_sem_up(void)
{
  int id;
  argint(0, &id);
  if (id < 0 || id >= NSEM)
    return -1;
  struct sem *s = &sems[id];

  acquire(&s->lock);
  s->value++;
  wakeup(s);
  release(&s->lock);
  return 0;
}
\end{lstlisting}

\begin{lstlisting}[caption={sem.o in the kernel and the three table entries}]
$ grep -n "sem.o" Makefile
21:  $K/sem.o \
$ grep -n "sys_sem" kernel/syscall.c
107:extern uint64 sys_sem_init(void);
108:extern uint64 sys_sem_down(void);
109:extern uint64 sys_sem_up(void);
138:  [SYS_sem_init] = sys_sem_init,
139:  [SYS_sem_down] = sys_sem_down,
140:  [SYS_sem_up]   = sys_sem_up,
\end{lstlisting}

\begin{lstlisting}[caption={Check 2: semtest}]
$ semtest
parent: done, I let the child go
child: I ran after the parent
ping 1
  pong 1
ping 2
  pong 2
ping 3
  pong 3
semtest: OK
\end{lstlisting}

The parent waits 10 ticks before it lets the child go, yet the child prints second: it was
asleep in \code{sem\_down}. Ping and pong alternate strictly, three times each. With
\code{trace} from Lab 10 the new calls can be watched directly (mask
\(2^{25} + 2^{26} = 100663296\)):

\begin{lstlisting}[caption={The semaphore calls of semtest, seen with trace}]
$ trace 100663296 semtest
parent: done, I let the child go
12: syscall sem_up -> 0
13: syscall sem_down -> 0
child: I ran after the parent
12: syscall sem_down -> 0
ping 1
12: syscall sem_up -> 0
14: syscall sem_down -> 0
  pong 1
14: syscall sem_up -> 0
12: syscall sem_down -> 0
ping 2
[... the same pattern for ping 2, pong 2, ping 3, pong 3 ...]
semtest: OK
\end{lstlisting}

The \code{sem\_down} of the child (pid 13) returns only after the \code{sem\_up} of the parent
(pid 12): the call was entered earlier, but it finished, and was printed, only when the child
was woken.

\section{Part 3. A semaphore as a lock: mix}

\code{mix} starts two processes that each print five lines, one letter per \code{write()}
system call. QEMU runs 3 CPUs, so the two processes really run at the same time.

\begin{lstlisting}[caption={mix without the lock}]
$ mix
AAAAbbbAbbAbAbAbAbAbAbAbAbbbAAbAAAAAbAbAbAbAbAbAAAbAbAb
bAbAbAbAbA
AAbAbAbAbAbAbbAbAbAbAbbbAAbAbAbAbAbAbbAbbAbbbbbbbb
bbAbAbAbAbAb
bAbAbAbAbAbAbbAbAbAAAAbAbAbAbbbAbAbAbAbAbAbAbA
AAAAbAbAbbAb
bbbbAbAbAbAbAbAbAbAbbAbAbAbAbAbAbAbAbAbAbAbAbA
AbAbAbAbAbAbAbAA
bAbAbAbbAbAbAbAbAbAbAbAbAbAbAbAbAbAbAbAbAbAbA
AAAAAAAA
\end{lstlisting}

\begin{lstlisting}[caption={mix with the lock}]
$ mix lock
AAAAAAAAAAAAAAAAAAAAAAAAAAAAAA
AAAAAAAAAAAAAAAAAAAAAAAAAAAAAA
AAAAAAAAAAAAAAAAAAAAAAAAAAAAAA
AAAAAAAAAAAAAAAAAAAAAAAAAAAAAA
AAAAAAAAAAAAAAAAAAAAAAAAAAAAAA
bbbbbbbbbbbbbbbbbbbbbbbbbbbbbb
bbbbbbbbbbbbbbbbbbbbbbbbbbbbbb
bbbbbbbbbbbbbbbbbbbbbbbbbbbbbb
bbbbbbbbbbbbbbbbbbbbbbbbbbbbbb
bbbbbbbbbbbbbbbbbbbbbbbbbbbbbb
\end{lstlisting}

\textbf{Observations:}
\begin{enumerate}
  \item \textbf{Why the letters mix.} One line costs 31 system calls: one \code{write()} for
  each of the 30 letters and one for the newline. Between any two of them the other process, on
  another CPU, can make its own \code{write()}, so the two lines are interleaved letter by
  letter.
  \item \textbf{The critical section.} \code{sem\_down(3);} before the inner loop and
  \code{sem\_up(3);} after it. Semaphore 3 starts at 1, so only one process at a time can be
  between the two calls, and a whole line is printed before the other process may start one.
  \item \textbf{Why \code{while} and not \code{if}.} \code{wakeup(s)} wakes every process that
  sleeps on \code{s}, but \code{sem\_up} added only one unit. The first process that gets the
  spinlock takes the unit and the value is 0 again; the others wake up, find 0 and must sleep
  again. A process that calls \code{sem\_down} without ever sleeping can also take the unit
  before the woken one gets the lock. With \code{if} the value would go below 0.
  \item \textbf{Why release before \code{sleep()}.} A spinlock is held with interrupts off, and
  whoever else wants it spins on a CPU. If the process slept while holding \code{s->lock},
  \code{sem\_up} could never acquire it to add 1 and call \code{wakeup}, so the sleeper would
  never be woken and the CPU that called \code{sem\_up} would spin forever: a deadlock inside
  the kernel. \code{sleep\_prepare} is called before the release so that a \code{sem\_up}
  arriving in between is not missed.
\end{enumerate}

The optional \code{sem\_trydown} call was not done.

\plainsection{Conclusions}

Building semaphores in the kernel showed how little they are made of. Each one is an integer
and a spinlock. Taking a unit is a loop: while the value is 0 the process records what it is
waiting for, releases the spinlock and sleeps, using no CPU; giving a unit back is an increment
and a wakeup. The user side is the recipe of Lab 10 repeated three times: a number, a
declaration, a stub and a table entry.

The two tests then reproduced Labs 11 and 12 on my own kernel. A semaphore that starts at 0
made the child wait for the parent, two semaphores passed the turn between ping and pong, and
a semaphore that starts at 1 turned a line of 31 separate system calls into a critical section,
so the output of two processes on three CPUs stopped mixing. The questions about \code{while}
and about releasing the spinlock before sleeping are the same lessons as the condition variable
and the deadlock of Lab 12, seen this time from inside the kernel.

\end{document}
